\section{Performance}\label{Performance}
\pagestyle{mypagestyle}
\begin{spacing}{1.75}
	
	\MUFGfont{
		\subsection{Performance objectives}
		Performance testing for FX simulation focuses on whether simulated distributions and dynamics are adequate for exposure generation:
		\begin{itemize}
			\item distributional fit of standardized innovations / PIT;
			\item quantile exceedance behaviour over rolling horizons;
			\item stability of calibrated volatility and sensitivity to configuration.
		\end{itemize}
	}
	
	\MUFGfont{
		\subsection{Backtesting design (recommended)}
		Let realized FX be $S_{t_i}$ sampled at the model frequency.
		Define log-return $r_i=\ln(S_{t_{i+1}}/S_{t_i})$.
		Under the model, standardized residuals should satisfy:
		\[
		\hat{z}_i=\frac{r_i-\hat{\mu}_i\Delta t}{\hat{\sigma}\sqrt{\Delta t}}\sim \mathcal{N}(0,1),
		\]
		with $\hat{\mu}_i$ implied by the forward drift (or set to zero for short horizons if consistent with implementation).
		Then apply:
		\begin{itemize}
			\item KS / AD / CvM tests on $\{\hat{z}_i\}$;
			\item PIT tests: $u_i=\Phi(\hat{z}_i)$ should be Uniform(0,1);
			\item exception tests for selected quantiles (e.g., 95\%, 99\%) over rolling windows;
			\item stability checks: re-estimate $\sigma$ on rolling windows and test for breaks.
		\end{itemize}
	}
	
	\MUFGfont{
		\subsection{Sensitivity tests}
		Mandatory sensitivities:
		\begin{itemize}
			\item window length: 1y vs 3y vs 5y;
			\item return frequency: 1-day vs 5-day;
			\item EWMA vs simple stdev;
			\item stressed scaling overlays (e.g., max of current vol and stress percentile).
		\end{itemize}
		Each sensitivity should be evaluated on impact to PFE metrics for representative portfolios (FX forwards, XCCY swaps, FX options if present).
	}
	
\end{spacing}
\clearpage
